% Source presentations for register row N-c-002.

\subsection{Separate zero--one-law presentations}
\label{sec:p09-presentation-n-c-002}

\subsubsection{Fagin's extension theory}

Let \(Y\) be a finite relational vocabulary whose relation symbols have
positive arity.  Fagin writes \(A_n(Y)\) for the \(Y\)-structures with universe
\(\{1,\ldots,n\}\) and, for a sentence
\(\varphi\), defines its labelled probability by
\[
 p_n(\varphi)
 =
 \frac{\lvert\{\mathcal A\in A_n(Y):\mathcal A\models\varphi\}\rvert}
      {\lvert A_n(Y)\rvert}.
\]
Equality is available, while constants and function symbols are absent from
the main theorem's vocabulary \cite[pp.~50--51]{Fagin1976}.

For pairwise distinct variables \(\bar x=(x_1,\ldots,x_m)\), a complete open
description \(M(\bar x)\) decides every equality and relational atom on that
tuple.  A compatible description \(N(\bar x,y)\) additionally decides every
atom involving a fresh point \(y\).  Fagin's theory \(T_Y\) contains, for every
such compatible pair, the extension sentence
\[
 \forall\bar x\,
 \bigl(M(\bar x)\to\exists y\,N(\bar x,y)\bigr).
\]
Here \(T_Y\) is a family of first-order sentences about \(Y\)-structures, not
a probability distribution and not a name for the countable random graph.
Its use is to turn an asymptotic counting assertion into a complete theory.

Fagin proves that \(T_Y\) is consistent and complete and that every member of
\(T_Y\) has limiting probability one.  The latter step is a finite estimate:
for any fixed requested one-point extension type, each remaining element has a
fixed positive probability of realizing it, and a union bound over the finitely
many parameter tuples makes the probability of any missing witness tend to
zero.  Completeness then gives the dichotomy
\[
 T_Y\models\varphi
 \quad\Longleftrightarrow\quad
 \lim_{n\to\infty}p_n(\varphi)=1;
\]
if \(T_Y\models\neg\varphi\), the limit is zero.  This is the source route to
the public zero--one law in Theorem~\ref{thm:p03-fol-zero-one}
\cite[thms.~1--4]{Fagin1976}.  The completeness proof uses back-and-forth for
countable models and the \L o\'s--Vaught test; the probability calculation is
displayed for a binary relation and then extended to finite relational
vocabularies.  It is not presented as a full modern proof with every general-
arity estimate expanded.

The article also changes the probability space.  If \(B_n(Y)\) contains one
representative of each isomorphism type and \(v_n(\varphi)\) is the proportion
of those representatives satisfying \(\varphi\), Theorem~9 proves that this
unlabelled limit exists and agrees with the labelled zero or one
\cite[pp.~54--55]{Fagin1976}.  The two measures are separate definitions; the
unlabelled result cannot be obtained by silently quotienting the displayed
formula for \(p_n\).

The vocabulary restrictions are witnessed by counterexamples.  If nullary
relation symbols were admitted, the sentence consisting of one such symbol
would already have probability \(1/2\).  With a unary predicate \(U\) and a
constant \(c\), the sentence \(U(c)\) likewise has limiting probability
\(1/2\).  With a unary function, the probability that there is no fixed point
tends to \(1/e\).  A second-order sentence need not have any limit at all.
Thus the positive-arity relational, constant-free, first-order hypotheses are
part of the theorem rather than stylistic simplifications
\cite[pp.~55--57]{Fagin1976}.

\subsubsection{Libkin's game and computable-random-graph route}

For a class \(\mathcal C\) of finite \(\sigma\)-structures, Libkin writes
\(\mu_n(\mathcal C)\) for the uniform proportion among structures with
universe \(\{0,\ldots,n-1\}\), and \(\mu(\mathcal C)\) for its limit when
that limit exists.  A logic has the zero--one law when every Boolean query it
defines has limiting probability zero or one
\cite[def.~12.1]{Libkin2004}.

For finite disjoint vertex sets \(U,V\) of sizes \(u,v\), the graph extension
sentence asks for a vertex outside \(U\cup V\) adjacent to every member of
\(U\) and to no member of \(V\).  In \(G(n,1/2)\), the probability that one
fixed pair \(U,V\) has no witness is
\[
 \bigl(1-2^{-(u+v)}\bigr)^{n-u-v}.
\]
There are at most \(n^{u+v}\) ordered choices of the two sets, so the
probability that any instance of this fixed extension requirement fails is at
most
\[
 n^{u+v}\bigl(1-2^{-(u+v)}\bigr)^{n-u-v},
\]
which tends to zero.  Libkin generalizes the complete-description schema from
graphs to every finite purely relational vocabulary
\cite[prop.~12.4 and pp.~240--241]{Libkin2004}.

The logical step is a game argument.  A sufficiently rich finite collection
of extension sentences gives Duplicator the required finite-pebble strategy;
formulas using only finitely many variables cannot distinguish structures
satisfying the corresponding collection.  In the source,
\(L^{\omega}_{\infty\omega}=\bigcup_{k\in\mathbb N}L^k_{\infty\omega}\)
is the infinitary language whose individual formulas use only finitely many
variables.  Theorem~12.2 states that this language has the zero--one law, and
Corollary~12.3 obtains the result for first-order, least fixed-point,
inflationary fixed-point, and partial fixed-point logic from their inclusion in
that finite-variable infinitary language
\cite[def.~11.1, thm.~11.2, thm.~12.2, and cor.~12.3]{Libkin2004}.

Libkin also supplies one computable countable model of the almost-sure graph
theory.  Write \(\operatorname{BIT}(i,j)\) when the \(j\)-th bit in the binary
expansion of \(i\) is \(1\).  Definition~12.8 puts an undirected edge between
\(i\) and \(j\), for \(j<i\), exactly when
\(\operatorname{BIT}(i,j)\) holds.  The larger vertex is therefore the first
argument of \(\operatorname{BIT}\), and the smaller vertex is the bit position.
The resulting graph satisfies every extension sentence.  Back-and-forth makes
the extension theory complete and \(\aleph_0\)-categorical; because the axioms
are recursive and complete, the theory is decidable
\cite[def.~12.8 and thm.~12.9]{Libkin2004}.  Corollaries~12.10--12.11 identify
truth in this graph with limiting probability one and make the almost-sure
theory decidable.  This explicit \(\operatorname{BIT}\) presentation belongs
to Libkin's account, not to Fagin's 1976 notation.

The presentation also marks its limit.  Adding a built-in order destroys these
zero--one results, and even successor makes the first-order law fail.  Even
cardinality has no zero--one truth limit, while connectivity can have limiting
probability one without being first-order definable
\cite[pp.~235--242]{Libkin2004}.  A zero--one law is therefore neither an
expressibility theorem nor an algorithm for finite validity.

\paragraph{Independent-priority boundary.}
The register separately names the 1969 paper of Glebskii, Kogan, Liogon'kii,
and Talanov.  Its identity and independent priority are established, but its
technical text was not available for source-specific verification.  The
probability space, notation, proof method, and examples of that paper are not
reconstructed by placing Fagin's complete-theory argument under the four
authors' names \cite{GlebskiiEtAl1969}.
